\subsection{Algebraic rank strata and Fitting ramification}
\label{sec:global-pencil-strata}

The pencil classification admits a global scheme-theoretic formulation.
On \(X^\times\), let \(\mathcal L=\OO_X(-1)\), and let
\(\mathcal A=P\mathcal L\), \(\mathcal B=Q\mathcal L\)
be its two image line bundles in \(E\otimes\OO_X\).  Set
\(\mathcal M=\mathcal A\oplus\mathcal B\).  The tautological
line \(\mathcal L\subset\mathcal M\) is the diagonal under
the isomorphisms induced by \(P\) and \(Q\).  Restricting the
Pluecker quadrics gives a morphism
\begin{equation}\label{eq:global-pencil-bundle-morphism}
 \mathcal I_{\mathrm{Pl},2}\otimes\OO_{X^\times}
 \xrightarrow{\ \rho\ }
 \mathcal K:=\ker\bigl(\Sym^2\mathcal M^*
                         \longrightarrow(\mathcal L^*)^{\otimes2}\bigr).
\end{equation}
The target has rank two.  Multiplication identifies it with
\((\mathcal M/\mathcal L)^*\otimes\mathcal M^*\).
The map lands in this kernel because each quadric vanishes on the
original Pluecker point.  A local generator of \(\mathcal L\)
identifies \(\rho\) with the family
\((b-a)\ell_{F,w}(a,b)\).

\begin{lemma}[Residual Cartier sections over the rank-one scheme]
\label{lem:residual-cartier-sections}
Endow \(T=D_1\setminus D_0\) with the determinantal locally
closed scheme structure, not its reduction.  On a bundle chart the
restricted Pluecker ideal on \(\Pj(\mathcal M_T)\) is generated
by
\[
            (b-a)(ca+db),\qquad (c,d)=\OO_T.
\]
Let \(T_{\mathrm{sep}}\subset T\) be the open where the residual
section avoids both endpoints and the diagonal.  In these coordinates
it is \(T\setminus V(cd(c+d))\).  The closed pencil intersection
over this open is the disjoint union of the effective Cartier sections
\[
                    s=[1:1],\qquad r=[d:-c].
\]
All these statements commute with arbitrary base change, including
to a nonreduced base.  The residual map
\(\tau:T_{\mathrm{sep}}\to T_{\mathrm{sep}}\) is a regular
involution and its fixed-point scheme is empty.
\end{lemma}
\begin{proof}
On \(D_1\setminus D_0\) some matrix entry of \(\rho\) is
invertible.  Its two-by-two minors vanish in the structure sheaf;
row and column operations over that sheaf therefore put the matrix
in rank-one normal form.  The coefficient rows of the residual
linear forms generate a line direct summand, represented by a
unimodular pair \((c,d)\).  Thus the displayed generator is an
ideal-sheaf equality over \(T\), not just a fibrewise equation.

A unimodular linear form on a projective-line bundle cuts out a
Cartier section over any base.  For example on an open where its
second coefficient is a unit, its equation on the relevant affine
chart is a unit multiple of a monic degree-one polynomial.  Such a
polynomial is a non-zero-divisor over an arbitrary coefficient ring.
This proves the Cartier assertions for \(b-a\) and \(ca+db\).
The two forms have determinant \(-(c+d)\).  Where \(c+d\)
is a unit their zero sections are disjoint, so their ideal sheaves
are comaximal.  Their product equals their intersection, and the
Chinese remainder isomorphism identifies the divisor with
\(s(T_{\mathrm{sep}})\amalg r(T_{\mathrm{sep}})\).
The additional units \(c,d\) exclude the two pure endpoints.
This construction and these unit conditions are unchanged by any
base change.

The section \(r\) lies in the Grassmannian as a scheme: substitution
in every restricted Pluecker equation gives zero.  It has both Schur
components nonzero.  At this point the component frames are changed
from \((Pw,Qw)\) to \((dPw,-cQw)\), an invertible diagonal
change on \(T_{\mathrm{sep}}\).  Writing \(a=da',b=-cb'\), its equation becomes
\[ cd(b'-a')(da'+cb'). \]
Thus the new residual coefficient pair is \((d,c)\), up to a unit;
its section \([c:-d]\) in the new frames represents
\(cd(Pw+Qw)\), which is exactly \(s\).
Consequently its residual coefficient matrix again has rank one as
a matrix over the base: the remaining linear factor is unimodular,
and at least one generating coefficient is a unit.  Thus \(r\)
factors through \(D_1\setminus D_0\), including its determinantal
scheme structure.  Its residual section is \(s\) and is disjoint
from it and from the endpoints, so the map factors through
\(T_{\mathrm{sep}}\).  Changing trivializations multiplies
these equations by units and does not change their sections; the
local morphisms therefore glue.  Division by the other Cartier
factor a second time gives \(\tau^2=\mathrm{id}\) as morphisms.
The disjointness of the sections gives an empty equalizer of
\(\tau\) and the identity, proving the fixed-point-scheme assertion.
\end{proof}


\begin{proposition}[The rank stratification and its finite part]
\label{prop:global-pencil-rank-strata}
Let \(D_1\) be the closed subscheme defined by the two-by-two
minors of \(\rho\), and \(D_0\) the one defined by its
entries.  The rank strata
\[
 X^\times\setminus D_1,\qquad D_1\setminus D_0,\qquad D_0
\]
have ranks two, one and zero, respectively.  On the rank-one stratum
the residual point is a regular section of \(\Pj(\mathcal M)\).
Deleting its endpoint and diagonal loci gives a regular fixed-point-free
involution.  The restriction of \(\pi\) to
\(X^\times\setminus D_0\) is quasi-finite.

Put \(Y=\Pj(E_{175})\times\Pj(E_{35})\) and
\(\mathcal R=V(\Fitt_0\Omega_{X^\times/Y})\).  We call
\(\mathcal R|_{X^\times\setminus D_0}\) the \emph{Fitting
ramification locus of the quasi-finite part}.  This convention
records the failure to be unramified, equivalently the nonvanishing
of the relative differential module \cite[Lemma~29.36.14]{StacksUnramified};
it does not assert a globally finite or flat double cover.
The same Fitting ideal is defined on all of \(X^\times\).
Its local formula is
\begin{equation}\label{eq:ramification-evaluation-ideal}
              (\ell_{F,w}(1,1):F\in\mathcal I_{\mathrm{Pl},2}).
\end{equation}
On the rank-one finite stratum it is exactly the double-point
subscheme, defined locally by \(c+d=0\) when
\(\ell=ca+db\).  It is empty on the rank-two stratum.  In
particular \(\pi\) is unramified on \(G_4^\circ\).
\end{proposition}
\begin{proof}
The determinantal definitions are invariant under changes of bundle
frames.  Where a rank-one matrix entry is invertible, its two-column
coefficient vectors are all multiples of one unimodular vector
\((c,d)\); its image is a line subbundle of \(\mathcal K\).
Dividing by the equation of the diagonal gives the line of residual
linear forms, whose zero is a regular section of
\(\Pj(\mathcal M)\).  This also works on the indicated
scheme structures, not merely on their closed points.  The residual
section lies in \(X\), because every restricted Pluecker quadric
vanishes on it.  The Cartier-section and involution assertions, with their
nonreduced scheme structures, follow from
Lemma~\ref{lem:residual-cartier-sections}.  On the complement
of \(D_0\), all geometric fibres have length at most two.  The
morphism is of finite type, hence quasi-finite there.

Let \(U\subset\Pj(E)\) be the open where both components
are nonzero.  The projection \(U\to Y\) is a
\(\mathbf G_m\)-bundle.  In a local trivialization its relative
differential module is generated by the torus differential.  The
conormal exact sequence for \(X^\times\subset U\) presents
\(\Omega_{X^\times/Y}\) as the quotient of this line bundle
by the relative derivatives of the Pluecker quadrics.  At a point
represented by \(w\), put \(a=1,b=1+t\).  The derivative of
\((b-a)\ell_{F,w}(a,b)\) in the vertical direction is
\(\ell_{F,w}(1,1)\).  This proves
\eqref{eq:ramification-evaluation-ideal}, including its scheme
structure.  At rank one one coefficient multiplier is a unit, so the
ideal restricts to \((c+d)\).  At rank two two independent linear
forms cannot both vanish at the diagonal point.  The last assertion
now follows from Theorem~\ref{thm:global-schur-recombination}.
\end{proof}

For clarity, the lengths in the ambient closed pencil and in a fibre
of \(\pi\) are different in the endpoint case:
\begin{center}
\begin{tabular}{llll}
\toprule
Rank & Type & Closed pencil intersection & Fibre after endpoint deletion\\
\midrule
2 & reduced point & length 1 & length 1\\
1 & distinct secant & two reduced points & length 2\\
1 & double point & length 2, nonreduced & length 2, nonreduced\\
1 & pure endpoint & two reduced points & length 1\\
0 & flag line & \(\Pj^1\) & \(\mathbf G_m\)\\
\bottomrule
\end{tabular}
\end{center}
The first row is the whole of \(G_4^\circ\), a nonempty smooth
open of \(\Gr(4,10)\) of dimension 24 and codimension zero in
that Grassmannian.  Every other row has empty intersection with
\(G_4^\circ\).  Thus there is no positive-codimension exceptional
stratum in this geometric family whose dimension or generic
singularity remains to be computed.  The determinantal schemes may
have further geometry on the binary-Jacobian boundary; the preceding
assertions do not assign an expected codimension to those ambient
schemes or make \(\pi\) into a global double cover.
